% !TeX root = main.tex
% Mereset counter huruf agar kembali ke A
\setcounter{huruf}{0}

% Judul Bab III (Pusat)
\begin{center}
	\textbf{BAB III} \\
	\textbf{PENUTUP}
\end{center}
\addcontentsline{toc}{section}{BAB III PENUTUP}

\vspace{0.5cm}

% Kesimpulan menggunakan \hurufsection
\hurufsection{Simpulan}{
	Berdasarkan pembahasan yang telah diuraikan pada bab sebelumnya, dapat disimpulkan bahwa implementasi teknologi informasi dalam dunia perpustakaan bukan lagi sekadar pilihan, melainkan sebuah kebutuhan. Otomasi perpustakaan terbukti mampu meningkatkan efisiensi kerja pustakawan dan memberikan kemudahan akses bagi para pemustaka.

	Meskipun terdapat tantangan dalam hal infrastruktur dan biaya, manfaat jangka panjang yang diperoleh dalam hal pelestarian dan penyebaran informasi jauh lebih besar bagi kemajuan ilmu pengetahuan.
}

% Saran menggunakan \hurufsection
\hurufsection{Saran}{
	Diharapkan pihak pengelola perpustakaan dapat terus melakukan pembaruan (\textit{upgrade}) terhadap sistem yang digunakan agar tetap relevan dengan perkembangan zaman. Selain itu, pelatihan bagi sumber daya manusia juga perlu ditingkatkan agar pemanfaatan teknologi informasi dapat berjalan secara optimal dan profesional.
}